\section{Time-Series-Native Agent Harness}
\label{sec:method}

We address the misalignment between language agents and time-series tasks through an \textit{agent harness}~\cite{ning2026code}: a time-series-native runtime that surrounds a frozen LLM policy and reshapes how it interacts with contextualized temporal data. Instead of forcing numerical signals into the token stream, the harness hosts temporal objects in a structured runtime environment and exposes them only through typed executable interfaces. The agent analyzes them through auditable temporal operations, while intermediate analytical states are maintained in a structured workspace.

\subsection{Harness Framework Overview}
\label{sec:overview}

\method instantiates this harness through three coordinated components. \textit{Executable temporal tools} (Section~\ref{sec:tools}) provide typed interfaces for accessing and analyzing time-series objects, producing grounded and auditable analytical evidence. \textit{Experience-driven capability evolution} (Section~\ref{sec:evolve}) abstracts recurring procedures as useful temporal routines to be reused rather than repeatedly rediscovered. \emph{Episodic multimodal memory} (Section~\ref{sec:memory}) stores past reasoning traces and retrieves relevant experience through complementary textual and temporal modalities.
Formally, \method augments a frozen LLM policy $\pi_\theta$ with a harness $\mathcal{H}$:
\begin{equation}
    \mathcal{H} = (\mathcal{T}, \mathcal{E}, \mathcal{M} ),
\end{equation}
where $\mathcal{T}$ is the set of executable temporal tools, $\mathcal{E}$ is the capability-evolution loop and $\mathcal{M}$ is the episodic multimodal memory.

\paragraph{Workspace and rollout.}
For each task $\tau=(q,\mathcal{X},\mathcal{Y},y^\star,\ell)$, $\mathbf{X}_{1:T}$ is loaded into a task-local \emph{workspace} $\mathcal{W}$ hosted by the harness runtime that
can be accessed through protocols (e.g., Model Context Protocol \citep{mcp2024}). This
workspace preserves the series at full numerical precision. The policy interacts with the harness through a
trajectory
\begin{equation}
    \zeta=(a_1,o_1,\ldots,a_K,o_K,\hat{y}),
\end{equation}
where each action $a_r\in\mathcal{A}$ denotes an agent decision, such as executing a runtime operation or performing a reasoning step. The corresponding observation $o_r\in\mathcal{O}$ records the outcome of this decision, either from the harness or the agent's internal reasoning. The policy selects
\begin{equation}
    a_r\sim\pi_\theta(\cdot\mid q,\mathcal{C},o_{1:r-1}),
\end{equation}
conditioned on the instruction, context, and accumulated evidence, ultimately emitting $\hat{y}\in\mathcal{Y}$.


\subsection{Runtime-Native Temporal Tools}
\label{sec:tools}

Despite differing in domains and objectives, time-series tasks often rely on a shared list of common temporal operations. Unlike conventional tool-augmented agents that must pass data through language messages, \method executes these operations inside the harness runtime, where the full time-series object already resides in the workspace $\mathcal{W}$. Thus, the agent only specifies the intended operation and its parameters without the need to spend context budget on tool interaction. Each tool runs directly on $\mathcal{W}$ at numerical precision and returns a compact structured observation. We initialize $\mathcal{T}$ with common time series tools and allow new tools to emerge through capability evolution, as summarized in Table~\ref{tab:tools}.

\paragraph{Auditable trajectories as reusable experience.}
The harness requires every numerical claim to be grounded in a returned tool observation rather than inferred from free-form language reasoning.
Each tool-related pair $(a_r,o_r)$ in the rollout also records the analytical intent behind that invocation. The resulting evidence chain ties the final answer $\hat{y}$ to verifiable temporal operations and forms a reusable problem-solving trace, which further supports capability evolution and episodic memory.



\subsection{Experience-Driven Capability Evolution}
\label{sec:evolve}

Runtime-native tools cover common temporal primitives, but real workflows often reveal recurring task-specific routines that cannot all be anticipated at design time.
The second facet of the harness turns such repeated experience into executable capability. Specifically, an evolution operator $\mathcal{E}$ identifies verified recurring sub-procedures and admits them into the tool set $\mathcal{T}$. As a result, $\mathcal{T}$ evolves and expands with the target task distribution.


\paragraph{Capability evolution mechanism.}
We implement the evolution operator $\mathcal{E}$ as a code-specialized LLM. Given clustered trajectories and task descriptions, $\mathcal{E}$ abstracts recurring analytical sub-procedures into self-contained executable tools, each with an invocation docstring and an implementation that operates on the runtime workspace $\mathcal{W}$. Each candidate tool is verified by execution on a held-out subset of similar tasks and is admitted into $\mathcal{T}$ only if its held-out success rate exceeds a threshold $\gamma$. Once admitted, the tool is immediately registered alongside the seed tools and can be invoked in future rollouts.




\subsection{Episodic Multimodal Memory}
\label{sec:memory}


\begin{table*}[t]
\caption{Overall results on the CiK benchmark. We report token usage when applicable. The best and second-best results in each metric column are highlighted in \textbf{bold} and \underline{underlined}, respectively. \method achieves the best average RCRPS and sMAPE, obtains the best RCRPS in 4 out of 5 context types, and is nearly tied regarding the remaining context type. Meanwhile, \method achieves substantially better performance with nearly half the token budget compared with the second-best multi-agent solution.}
\label{table:main-results}
\centering
\resizebox{\textwidth}{!}{%
\begin{tabular}{lcccccccc}
\toprule
 \multirow{2}{*}{\makecell{Method}} 
 & \multirow{2}{*}{\makecell{Average \\ RCRPS ($\downarrow$)}} 
 & \multirow{2}{*}{\makecell{Average \\ sMAPE ($\downarrow$)}} 
 & \multirow{2}{*}{\makecell{Average \\ Token Usage ($\downarrow$)}}
 & \multicolumn{5}{c}{RCRPS by Context Type ($\downarrow$)} \\
\cmidrule(lr){5-9}
 & & & & Intemporal & Historical & Future & Covariate & Causal \\
\midrule
\multicolumn{9}{l}{\textit{Traditional Time Series Models}} \\
~~~ARIMA \citeyearpar{broomhead1989time} & 0.2772 & 91.2514 & -- & 0.2987 & 0.1032 & 0.1577 & 0.2073 & 0.2822 \\
~~~ETS \citeyearpar{jain2017study} & 0.3282 & 84.5230 & -- & 0.3654 & 0.1738 & 0.1966 & 0.2402 & 0.3290 \\
% \midrule
% \multicolumn{9}{l}{\textit{Time Series Neural Networks}} \\
~~~DLinear \citeyearpar{DBLP:conf/aaai/ZengCZ023} & 2.3448 & 184.1755 & -- & 3.0988 & 4.9004 & 2.6461 & 1.9710 & 3.0394 \\
~~~PatchTST \citeyearpar{DBLP:conf/iclr/NieNSK23} & 2.3374 & 166.8354 & -- & 3.0980 & 4.9007 & 2.6294 & 1.9614 & 3.0523 \\
\midrule
\multicolumn{9}{l}{\textit{Time Series Foundation Models}} \\
~~~Chronos-2 \citeyearpar{DBLP:journals/corr/abs-2510-15821} & 0.2344 & 90.8453 & -- & 0.2277 & 0.1852 & 0.2262 & 0.2406 & 0.3426 \\
~~~Lag-Llama \citeyearpar{rasul2023lag} & 0.2628 & 95.1976 & -- & 0.2475 & 0.1233 & 0.1987 & 0.2461 & 0.3299 \\
~~~Moirai-Large \citeyear{DBLP:conf/icml/WooLKXSS24} & 0.4545 & 107.8152 & -- & 0.3865 & 0.0929 & 0.3936 & 0.3792 & 0.2690 \\
\midrule
\multicolumn{9}{l}{\textit{LLM/Agent on Time Series}} \\
~~~UniTime \citeyearpar{DBLP:conf/www/LiuHLDLHZ24} & 0.2822 & 93.3854 & -- & 0.2910 & \underline{0.0889} & 0.1665 & 0.2964 & 0.2677 \\
~~~Time-LLM \citeyearpar{DBLP:conf/iclr/0005WMCZSCLLPW24} & 0.3990 & 102.6825 & -- & 0.3629 & 0.1144 & 0.3172 & 0.3589 & 0.3146 \\
~~~TS-Agent \citeyearpar{DBLP:journals/corr/abs-2510-07432} & 0.1421 & 61.2980 & 47455.5880 & 0.1434 & 0.0904 & 0.1210 & 0.1516 & \textbf{0.1596} \\
~~~TSci \citeyearpar{DBLP:journals/corr/abs-2510-01538} & 0.1448 & 69.9068 & 47905.1056 & 0.1540 & 0.1001 & 0.1181 & 0.1694 & 0.1971 \\
\midrule
\multicolumn{9}{l}{\textit{General Agentic Pipelines}} \\
~~~Direct Prompt & 0.1703 & 60.4407 & 24861.3451 & 0.1938 & 0.1244 & 0.1239 & 0.1880 & 0.2793 \\
~~~CoT Prompt \citeyearpar{DBLP:conf/nips/Wei0SBIXCLZ22} & 0.1726 & 66.3724 & 27120.8662 & 0.1771 & 0.1230 & 0.1520 & 0.1932 & 0.2392 \\
~~~ReAct \citeyearpar{DBLP:conf/iclr/YaoZYDSN023} & 0.1514 & 66.2284 & 25502.9507 & 0.1678 & 0.1452 & 0.1317 & 0.1681 & 0.2249 \\
~~~Self-Reflection \citeyearpar{DBLP:journals/corr/abs-2405-06682} & 0.1608 & 64.7013 & 28942.8662 & 0.1869 & 0.1240 & 0.1131 & 0.1851 & 0.2756 \\
~~~Multi-Agent Reflection \citeyearpar{DBLP:conf/icml/YuanX25} & \underline{0.1294} & \underline{55.5704} & 63033.0070 & \underline{0.1346} & 0.1086 & \underline{0.1108} & \underline{0.1437} & 0.1697 \\
\midrule
\method (Ours) & \textbf{0.1145} & \textbf{52.4542} & 35553.2113 & \textbf{0.1303} & \textbf{0.0886} & \textbf{0.0703} & \textbf{0.1239} & \underline{0.1597} \\
\bottomrule
\end{tabular}
}
\vspace*{-0.2cm}
\end{table*}



Episodic memory is common in agentic systems, but its design is nontrivial in contextualized time-series settings, since relevance depends on both task framing and temporal structure. A text-based or jointly serialized retrieval space can (1) retrieve misleading precedents, since token similarity does not necessarily preserve numerical proximity (Proposition~\ref{prop:token_numeric}); (2) risk modality imbalance for long series; (3) exhibit scale instability, as normalization choices can change which temporal patterns appear similar.

\paragraph{Memory records.}
Each memory record $m\in\mathcal{M}$ stores a past task, its retrieval keys, and its successful analytical trace:
\begin{equation}
    m = (\tau, \phi, \psi, \zeta, \rho),
\end{equation}
where $\tau$ is the original task, $\zeta$ is the rollout trajectory produced under the harness, and $\rho$ is a short transferable rationale summarizing the contextual cue and analytical rule behind the solution. For example, in the energy domain, $\rho$ may note that a snowstorm scenario implies a sharp electricity demand increase due to air-conditioning load. The remaining fields $\phi$ and $\psi$ serve as retrieval keys over the signal and text modalities, respectively.


\paragraph{Multimodal retrieval keys.}
For holistic memory retrieval, we design retrieval keys for both modalities. The context key $\phi=\Phi(c_\tau)$ embeds a textual descriptor $c_\tau$ constructed from the task context. The time series key $\psi=\Psi(\mathbf{X}_{1:T})$ is a time-series fingerprint that encodes structural, statistical, spectral, temporal-dynamic, and multivariate properties. 
Due to page limitations, we provide details of fingerprint computation in Appendix \ref{ap:fingerprint}.


\paragraph{Similarity computation and retrieval.}
Given a query task $\tau_q$ with keys $(\psi_q,\phi_q)$, the text-side similarity is defined by cosine similarity,
\begin{equation}
    s_{\mathrm{text}}(q,m)=\cos(\phi_q,\phi_m),
\end{equation}
while the time series similarity is defined by the negative distance between normalized fingerprints,
\begin{equation}
    s_{\mathrm{ts}}(q,m)=-\|\tilde{\psi}_q-\tilde{\psi}_m\|_2,
\end{equation}
where $\tilde{\psi}$ denotes the normalized fingerprint. These two scores can be merged in different ways, such as weighted aggregation, reciprocal-rank fusion, or staged retrieval. In \method, we adopt a simple yet effective two-stage top-k retrieval strategy with details provided in Appendix \ref{ap:multimodal_memory}.